\documentclass[10pt,a4paper,twocolumn]{article}
\usepackage[utf8]{inputenc}
\usepackage[portuguese]{babel}
\usepackage[top=2cm, bottom=2cm, left=1.5cm, right=1.5cm, columnsep=0.6cm]{geometry}
\usepackage{amsmath,amssymb,amsfonts}
\usepackage{graphicx}
\usepackage{booktabs}
\usepackage{multirow}
\usepackage{url}
\usepackage[hidelinks]{hyperref}
\usepackage{float}
\usepackage{algorithm}
\usepackage{algpseudocode}

% ABNT para citações - substitui natbib
\usepackage[alf,abnt-emphasize=bf,abnt-and-type=&,abnt-etal-list=0,abnt-etal-text=it]{abntex2cite}

% Formatação de título
\usepackage{titling}
\pretitle{\begin{center}\LARGE\bfseries}
\posttitle{\par\end{center}\vskip 0.5em}
\preauthor{\begin{center}\normalsize}
\postauthor{\end{center}}
\predate{\begin{center}\small}
\postdate{\par\end{center}\vskip 1em}

% Configurações AJUSTADAS para eliminar espaços extras entre floats
\setlength{\textfloatsep}{8pt plus 2pt minus 2pt}
\setlength{\floatsep}{8pt plus 2pt minus 2pt}
\setlength{\intextsep}{8pt plus 2pt minus 2pt}
\setlength{\dbltextfloatsep}{8pt plus 2pt minus 2pt}

% Controle de float pages
\renewcommand{\floatpagefraction}{0.8}
\renewcommand{\topfraction}{0.9}
\renewcommand{\bottomfraction}{0.8}
\renewcommand{\textfraction}{0.1}

\addto\captionsportuguese{
  \renewcommand{\tablename}{Table}
}
\setlength{\belowcaptionskip}{4pt} 

\title{Evaluation of Algorithms for Fractal Dimension Estimation in Medical Images: Complexity and Performance Analysis}

\author{
Gabriel Martins Brum%\textsuperscript{1*}, Nome2\textsuperscript{2}, Nome3\textsuperscript{3}
\\
\normalsize g.brum@unesp.br\\
}

\date{}

\begin{document}

\twocolumn[
\maketitle

\begin{@twocolumnfalse}
\begin{abstract}
\noindent Este estudo avalia quatro algoritmos de estimação de dimensão fractal — Box-Counting Clássico (BC), Gliding Box (GB), Differential Box-Counting (DBC) e Blanket Method (BM) — sob a perspectiva da engenharia algorítmica aplicada a imagens médicas. Limites assintóticos formais de tempo e espaço auxiliar são derivados para formulações ingênuas e otimizadas, com validação empírica em um conjunto de dados histológicos de câncer de mama. Fundamentalmente, os resultados demonstram que melhorias assintóticas teóricas nem sempre se traduzem em ganhos práticos de tempo de execução. Embora otimizações estruturais tenham proporcionado acelerações substanciais para o GB (imagens integrais) e DBC (tabelas esparsas), elas introduziram severas sobrecargas de memória ou foram neutralizadas por ineficiências de cache e condições de parada antecipada, como observado no desempenho paradoxalmente inferior do BC otimizado. A análise de efetividade revela que, enquanto BC e BM fornecem forte separação estatística entre tecidos benignos e malignos ($p < 0.01$), o poder discriminatório do GB reside na lacunaridade em vez da dimensão fractal. Por fim, este trabalho expõe o trade-off crítico entre a eficiência computacional de métodos binários, inerentemente sensíveis a limiares, e a robustez dos métodos em tons de cinza. Esses achados fornecem um arcabouço rigoroso para a seleção de algoritmos em pipelines clínicos, enfatizando que a hierarquia de memória e restrições de implementação são componentes tão críticos para o desempenho em mundo real quanto a complexidade assintótica.

\vspace{0.3em}
\noindent \textbf{Palavras-chave:} Dimensão fractal; Box-counting; Complexidade algorítmica; Imagens médicas; Classificação histológica.
\end{abstract}

\renewcommand{\abstractname}{Abstract}
\begin{abstract}
\noindent This study evaluates four fractal dimension estimation algorithms—Classical Box-Counting (BC), Gliding Box (GB), Differential Box-Counting (DBC), and the Blanket Method (BM)—from an algorithmic engineering perspective applied to medical imaging. Formal asymptotic bounds for time and auxiliary space are derived across naive and structurally optimized formulations, with empirical validation on a breast cancer histology dataset. Crucially, our findings demonstrate that theoretical asymptotic improvements do not strictly translate into practical runtime gains. While structural optimizations yielded substantial speedups for GB (integral images) and DBC (sparse tables), they introduced severe memory overheads or were counteracted by cache inefficiencies and early-exit conditions, as observed in the paradoxically slower optimized BC. Effectiveness analysis reveals that while BC and BM provide strong statistical separation between benign and malignant tissues ($p < 0.01$), GB's discriminatory power relies on lacunarity rather than fractal dimension. Ultimately, this work exposes a fundamental trade-off between the computational efficiency of threshold-sensitive binary methods and the robust, albeit resource-intensive, nature of grayscale methods. These insights provide a rigorous framework for algorithm selection in clinical pipelines, emphasizing that memory hierarchy and dataset-specific implementation constraints are as critical to real-world performance as asymptotic complexity.

\vspace{0.3em}
\noindent \textbf{Keywords:} Fractal dimension; Box-counting; Algorithmic complexity; Medical imaging; Histological classification.
\end{abstract}

\end{@twocolumnfalse}

\vspace{1em}
]

% ============================================================
%  SECTION 1 — INTRODUCTION
% ============================================================
\section{Introduction}

Fractal geometry provides a mathematical framework for quantifying the structural complexity and statistical self-similarity of biological phenomena that classical Euclidean geometry fails to capture \cite{Mandelbrot1982}. In medical image analysis, the fractal dimension ($D$) has been widely adopted to characterize tissue morphology, demonstrating significant clinical utility in assessing pulmonary emphysema \cite{Mishima1999}, retinal vasculature \cite{Jelinek2000}, cortical folding \cite{Free1996}, and tumour margin irregularity \cite{Rangayyan2007}. Recently, fractal descriptors have evolved beyond standalone features, being integrated into hybrid machine learning pipelines to enhance model interpretability \cite{roberto2021fractal} and guide data augmentation strategies \cite{borgue2025association, Lopes2011}.

Despite this proven discriminative utility, existing literature focuses heavily on diagnostic effectiveness while largely neglecting the underlying computational trade-offs. Contemporary clinical datasets routinely comprise high-resolution images where algorithmic efficiency is no longer a secondary concern, but a strict practical constraint. Consequently, the critical research question must shift from \textit{whether} fractal dimension is discriminative to \textit{how efficiently and robustly} it can be extracted under realistic resource limitations.

This study investigates the relationship between theoretical algorithmic complexity, implementation choices, memory utilization, and practical execution time. We evaluate four fundamental estimation algorithms that span the principal design axes of fractal image analysis (binary versus grayscale inputs, and disjoint versus dense-overlapping covering strategies): Classical Box-Counting (BC) \cite{Mandelbrot1982}, Gliding Box (GB) \cite{Allain1991}, Differential Box-Counting (DBC) \cite{Sarkar1994}, and the Blanket Method (BM) \cite{Peleg1984}. 

We hypothesize that theoretical asymptotic improvements do not necessarily translate into practical runtime gains. While structural optimizations (such as integral images or sparse tables) reduce mathematical complexity, their real-world performance is heavily mediated by memory hierarchy bottlenecks, cache behavior, early-exit conditions, and preprocessing requirements. To test this hypothesis, this work addresses the following concise objectives:

\begin{itemize}
  \item \textbf{Theoretical vs. Empirical Scaling:} Derive formal asymptotic time and auxiliary space bounds for naive and structurally optimized formulations of all four methods, explicitly contrasting theoretical predictions with observed wall-clock performance on a standard medical dataset.
  \item \textbf{Resource Trade-offs:} Quantify the execution and memory overheads imposed by advanced data structures and threshold-based preprocessing steps, evaluating their feasibility for clinical deployment.
  \item \textbf{Discriminatory Effectiveness:} Assess the statistical separation between pathological and healthy tissues generated by each algorithm, culminating in a synthesized trade-off analysis that balances computational cost against diagnostic value.
\end{itemize}

The remainder of this article is structured as follows: Section~2 establishes the theoretical foundations of the evaluated methods. Section~3 details the algorithm designs and their asymptotic complexities. Section~4 outlines the experimental methodology. Section~5 analyzes the empirical results, explicitly addressing the divergence between theory and practice. Section~6 discusses threats to validity, and Section~7 concludes with evidence-driven guidelines for algorithm selection in medical image analysis.

% ============================================================
%  SECTION 2 — THEORETICAL FOUNDATION
% ============================================================
\section{Theoretical Foundation}
\label{sec:theory}

% ------------------------------------------------------------
\subsection{Fractal Analysis and Dimension Estimation}
\label{sec:theory:fractal}

Fractal sets exhibit \textit{statistical self-similarity}, maintaining structural complexity across observation scales~\cite{Mandelbrot1982}. This property is characterized by power-law scaling, where a characteristic measure $M(\varepsilon)$ obtained at scale $\varepsilon$ satisfies

\begin{equation}
  M(\varepsilon) \;\sim\; C \cdot \varepsilon^{-D},
  \label{eq:powerlaw}
\end{equation}

\noindent with $C$ acting as a proportionality constant and $D$ representing the fractal dimension~\cite{Falconer2003}. Taking logarithms of both sides yields a linear relationship:

\begin{equation}
  \log M(\varepsilon) = \log C - D \log \varepsilon.
  \label{eq:loglog}
\end{equation}

\noindent $D$ is estimated as the negative slope of a least-squares regression over the pairs $\{(\log \varepsilon_i, \log M(\varepsilon_i))\}_{i=1}^{m}$ across $m$ discrete scales.

As the theoretical Hausdorff dimension~\cite{Hausdorff1919} is intractable on discrete grids, the Minkowski--Bouligand (box-counting) dimension $D_B$ serves as the standard approximation~\cite{Falconer2003}:

\begin{equation}
  D_B = \lim_{\varepsilon \to 0}
        \frac{\log N(\varepsilon)}{\log(1/\varepsilon)},
  \label{eq:boxdim}
\end{equation}

\noindent where $N(\varepsilon)$ is the minimum number of boxes of side $\varepsilon$ needed to cover the set. The regression's coefficient of determination ($R^{2}$) validates the presence of consistent power-law scaling, a prerequisite for interpreting $D_B$ as a meaningful structural descriptor~\cite{Lopes2011}.

% ------------------------------------------------------------
\subsection{Classical Box-Counting Method}
\label{sec:theory:bc}

The Classical Box-Counting (BC) method discretizes Eq.~\eqref{eq:boxdim} for binary images~\cite{Mandelbrot1982,Liebovitch1989}. Given an image $I$ of $N = W \times H$ pixels, the domain is partitioned into $\lceil W/\varepsilon \rceil \times \lceil H/\varepsilon \rceil$ disjoint cells of side $\varepsilon$. A cell is considered \textit{occupied} if it contains at least one foreground pixel:

\begin{equation}
  N(\varepsilon)
    = \sum_{i=0}^{\lceil W/\varepsilon \rceil - 1}
      \sum_{j=0}^{\lceil H/\varepsilon \rceil - 1}
      \mathbf{1}\!\left[
        \exists\,(x,y) \in B_{ij}(\varepsilon) : I(x,y) = 1
      \right].
  \label{eq:bc_count}
\end{equation}

Applying BC to grayscale medical images requires prior binarization, typically via Otsu's global thresholding~\cite{Otsu1979}. While binarization reduces the representational complexity, it irreversibly discards intensity gradients, making $D_B$ highly sensitive to the threshold parameter $t^{*}$.

To bypass exhaustive pixel scanning, optimized implementations pre-compute the \textit{integral image} $P[x][y] = \sum_{i \leq x, j \leq y} I[i][j]$ via the recurrence~\cite{Crow1984}:

\begin{equation}
  P[x][y]
    = I[x][y] + P[x{-}1][y] + P[x][y{-}1] - P[x{-}1][y{-}1].
  \label{eq:prefix}
\end{equation}

\noindent The foreground pixel count within any rectangular cell is then extracted in $\mathcal{O}(1)$ time:

\begin{equation}
  \sigma(x_1,y_1,x_2,y_2) =
    \begin{aligned}[t]
      &P[x_2][y_2] - P[x_1{-}1][y_2] \\
      &- P[x_2][y_1{-}1] + P[x_1{-}1][y_1{-}1].
    \end{aligned}
  \label{eq:rect_sum}
\end{equation}

% ------------------------------------------------------------
\subsection{Gliding Box Method}
\label{sec:theory:gb}

The Gliding Box (GB) method~\cite{Allain1991} evaluates a dense sliding window rather than a disjoint grid. A box of side $L$ is centered at every valid pixel position $(x, y)$, defining the \textit{box mass} as:

\begin{equation}
  m_{x,y}(L) = \sum_{i=0}^{L-1} \sum_{j=0}^{L-1} I[x+i][y+j].
  \label{eq:gb_mass}
\end{equation}

\noindent Across $T(L) = (W-L+1)(H-L+1)$ valid placements, the \textit{mass-frequency distribution} $n(m, L)$ dictates the normalized probability $P(m, L) = n(m, L) / T(L)$. $D_B$ is estimated by substituting the mean box mass into Eq.~\eqref{eq:loglog}~\cite{Plotnick1996}:

\begin{equation}
  \mu(L) = \sum_{m=0}^{L^{2}} m \, P(m, L).
  \label{eq:gb_mu1}
\end{equation}

GB also extracts \textit{lacunarity}, capturing the translational heterogeneity of the spatial mass distribution:

\begin{equation}
  \Lambda(L)
    = \frac{\mu^{(2)}(L) - [\mu(L)]^{2}}{[\mu(L)]^{2}},
  \label{eq:lacunarity}
\end{equation}

\noindent where $\mu^{(2)}(L) = \sum_{m} m^{2} P(m, L)$. Homogeneous textures yield $\Lambda \approx 0$, whereas clustered geometries produce large $\Lambda$~\cite{Allain1991,Tolle2008}. Integral images (Eq.~\eqref{eq:rect_sum}) optimize mass evaluations to $\mathcal{O}(1)$~\cite{Crow1984}, a critical reduction given the dense $\Theta(N)$ overlapping evaluations per scale.

% ------------------------------------------------------------
\subsection{Differential Box-Counting Method}
\label{sec:theory:dbc}

The Differential Box-Counting (DBC) method~\cite{Sarkar1994} avoids binarization by modeling the grayscale image $I \in \{0, \ldots, G{-}1\}$ as a three-dimensional topographic relief. For a cell of side $\varepsilon$, the intensity axis is partitioned into discrete bands of height $\delta = \lceil G / \lceil W/\varepsilon \rceil \rceil$. The number of boxes required to encompass the intensity gradient within cell $(i, j)$ is:

\begin{equation}
  n_{ij}(\varepsilon)
    = \left\lceil \frac{I_{\max}^{ij} - I_{\min}^{ij}}{\delta}
      \right\rceil + 1,
  \label{eq:dbc_count}
\end{equation}

\noindent where $I_{\max}^{ij}$ and $I_{\min}^{ij}$ denote local extrema. The total covering count $N(\varepsilon) = \sum_{i,j} n_{ij}(\varepsilon)$ determines $D_B$. While DBC preserves the structural fidelity of grayscale inputs, optimized execution relies on 2D sparse tables to reduce the local extrema search from $\Theta(\varepsilon^2)$ to $\mathcal{O}(1)$ queries.

% ------------------------------------------------------------
\subsection{Blanket Method}
\label{sec:theory:bm}

The Blanket Method (BM)~\cite{Peleg1984} computes $D_B$ by quantifying the apparent area growth of the intensity surface under sequential morphological operations. The upper ($u_\varepsilon$) and lower ($b_\varepsilon$) \textit{blankets} at dilation level $\varepsilon$ are recursively defined:

\begin{align}
  u_\varepsilon(x,y) &= \max\!\Bigl(u_{\varepsilon-1}(x,y)+1,\;
    \max_{(s,t)\in\mathcal{N}} u_{\varepsilon-1}(s,t)\Bigr),
  \label{eq:blanket_u} \\
  b_\varepsilon(x,y) &= \min\!\Bigl(b_{\varepsilon-1}(x,y)-1,\;
    \min_{(s,t)\in\mathcal{N}} b_{\varepsilon-1}(s,t)\Bigr),
  \label{eq:blanket_b}
\end{align}

\noindent initialized at $u_0 = b_0 = I$, with $\mathcal{N}$ representing the 4-connected neighborhood. The enclosed \textit{blanket volume} is:

\begin{equation}
  V(\varepsilon)
    = \sum_{x,y} \bigl[u_\varepsilon(x,y) - b_\varepsilon(x,y)\bigr].
  \label{eq:blanket_vol}
\end{equation}

\noindent The apparent surface area $A(\varepsilon) = \bigl[V(\varepsilon) - V(\varepsilon{-}1)\bigr] / 2$ substitutes $M(\varepsilon)$ in Eq.~\eqref{eq:loglog} to derive the fractal dimension. Operating strictly on intensity profiles, BM avoids thresholding artifacts. While a naive implementation incurs a neighborhood-dependent computational cost, separable monotonic deque filters~\cite{Lemire2006} optimize the 2D morphological extrema to $\Theta(N)$ total operations per scale.

% ============================================================
%  SECTION 3 — ALGORITHM PROJECT AND ANALYSIS
% ============================================================
\section{Algorithm Project and Analysis}
\label{sec:algorithms}

This section details the algorithmic design and formal asymptotic analysis of the four fractal dimension estimators. For each method, we first define a naive brute-force formulation, identify its computational bottleneck, and subsequently derive a structurally optimized counterpart. Throughout this analysis, $N = W \times H$ denotes the total image pixel count, $m = |\mathcal{S}|$ the number of evaluated scales, and $\varepsilon$ (or $L$) the current box side length. Asymptotic bounds follow standard Bachmann–Landau notation ($\mathcal{O}$, $\Theta$, $\Omega$).

% ------------------------------------------------------------
\subsection{Classical Box-Counting (BC)}
\label{subsec:bc}
% ------------------------------------------------------------

\subsubsection{Naive Formulation}

Algorithm~\ref{alg:bc_naive} implements Eq.~\eqref{eq:bc_count} through exhaustive cell scanning. For each scale $\varepsilon$, the grid contains $\Theta(N/\varepsilon^{2})$ cells. In the worst case (a fully foreground cell), evaluating a cell requires $\varepsilon^{2}$ pixel inspections; in the best case (first pixel is foreground), it requires one. The worst-case per-scale cost is:

\begin{equation}
  T_{\text{BC,naive}}(\varepsilon)
    = \Theta\!\left(\frac{N}{\varepsilon^{2}}\right)
      \cdot \mathcal{O}(\varepsilon^{2})
    = \mathcal{O}(N).
\end{equation}

The strict lower bound is $\Omega(N/\varepsilon^{2})$ since every cell must be visited. Summing over all $m$ scales, including the $\Theta(N)$ binarization pass, yields the total execution time:

\begin{equation}
  T_{\text{BC,naive}}(N,m)
    = \Theta(N) + \sum_{\varepsilon \in \mathcal{S}}
      \mathcal{O}(N)
    = \mathcal{O}(mN).
  \label{eq:bc_naive_time}
\end{equation}

Beyond the read-only input buffer, auxiliary space is $\Theta(1)$.

\begin{algorithm}[t]
\caption{Naive Classical Box-Counting}
\label{alg:bc_naive}
\begin{algorithmic}[1]
\Require Binary image $I[0\ldots W{-}1][0\ldots H{-}1]$, scale set $\mathcal{S}$
\Ensure Pairs $\{(\varepsilon,\, N(\varepsilon)) : \varepsilon \in \mathcal{S}\}$
\State Binarize $I$ via Otsu's threshold $t^{*}$ \cite{Otsu1979}
\For{each $\varepsilon \in \mathcal{S}$}
  \State $\mathit{count} \leftarrow 0$
  \For{$b_x \leftarrow 0$ \textbf{to} $\lceil W/\varepsilon \rceil - 1$}
    \For{$b_y \leftarrow 0$ \textbf{to} $\lceil H/\varepsilon \rceil - 1$}
      \State $\mathit{occupied} \leftarrow \textsc{False}$
      \For{$dx \leftarrow 0$ \textbf{to} $\varepsilon - 1$}
        \For{$dy \leftarrow 0$ \textbf{to} $\varepsilon - 1$}
          \State $px \leftarrow b_x \cdot \varepsilon + dx$;\; $py \leftarrow b_y \cdot \varepsilon + dy$
          \If{$px < W \;\wedge\; py < H \;\wedge\; I[px][py] = 1$}
            \State $\mathit{occupied} \leftarrow \textsc{True}$;\; \textbf{break}
          \EndIf
        \EndFor
      \EndFor
      \If{$\mathit{occupied}$}
        \State $\mathit{count} \leftarrow \mathit{count} + 1$
      \EndIf
    \EndFor
  \EndFor
  \State \textbf{yield} $(\varepsilon,\, \mathit{count})$
\EndFor
\State Estimate $D_B$ by least-squares regression on $\{(\log \varepsilon_i,\, \log N(\varepsilon_i))\}$
\end{algorithmic}
\end{algorithm}

% ............................................................
\subsubsection{Optimized Formulation}

The inner double loop of Algorithm~\ref{alg:bc_naive} represents the primary computational bottleneck. Algorithm~\ref{alg:bc_opt} eliminates this by pre-computing the integral image $P$ (Eq.~\eqref{eq:prefix}). Consequently, region sums are resolved in $\mathcal{O}(1)$ via Eq.~\eqref{eq:rect_sum}~\cite{Crow1984}, determining cell occupancy ($\sigma > 0$) without pixel scanning.

\begin{algorithm}[t]
\caption{Optimized Box-Counting (Integral Image)}
\label{alg:bc_opt}
\begin{algorithmic}[1]
\Require Binary image $I[0\ldots W{-}1][0\ldots H{-}1]$, scale set $\mathcal{S}$
\Ensure Pairs $\{(\varepsilon,\, N(\varepsilon)) : \varepsilon \in \mathcal{S}\}$
\State Binarize $I$ via Otsu's threshold $t^{*}$ \cite{Otsu1979}
\State Build integral image $P$ via Eq.~\eqref{eq:prefix} \hfill \Comment{$\Theta(N)$ time and space \cite{Crow1984}}
\For{each $\varepsilon \in \mathcal{S}$}
  \State $\mathit{count} \leftarrow 0$
  \For{$b_x \leftarrow 0$ \textbf{to} $\lceil W/\varepsilon \rceil - 1$}
    \For{$b_y \leftarrow 0$ \textbf{to} $\lceil H/\varepsilon \rceil - 1$}
      \State $x_1 \leftarrow b_x \cdot \varepsilon$;\; $y_1 \leftarrow b_y \cdot \varepsilon$
      \State $x_2 \leftarrow \min(x_1{+}\varepsilon{-}1, W{-}1)$;\; $y_2 \leftarrow \min(y_1{+}\varepsilon{-}1, H{-}1)$
      \If{$\sigma(x_1, y_1, x_2, y_2) > 0$} \hfill\Comment{Eq.~\eqref{eq:rect_sum}, $\mathcal{O}(1)$}
        \State $\mathit{count} \leftarrow \mathit{count} + 1$
      \EndIf
    \EndFor
  \EndFor
  \State \textbf{yield} $(\varepsilon,\, \mathit{count})$
\EndFor
\State Estimate $D_B$ by least-squares regression on $\{(\log \varepsilon_i,\, \log N(\varepsilon_i))\}$
\end{algorithmic}
\end{algorithm}

The prefix sum is computed once in $\Theta(N)$. Subsequently, traversing the $\Theta(N/\varepsilon^{2})$ cells per scale costs $\mathcal{O}(1)$ each:

\begin{equation}
  T_{\text{BC,opt}}(N, m)
    = \Theta(N)
    + \sum_{\varepsilon \in \mathcal{S}}
      \Theta\!\left(\frac{N}{\varepsilon^{2}}\right).
  \label{eq:bc_opt_time}
\end{equation}

For a geometric scale set (e.g., $\varepsilon \in \{2, 4, 8 \ldots\}$), the sum is dominated by the smallest scale, collapsing the total time complexity to $\Theta(N)$. This strictly improves upon the $\mathcal{O}(mN)$ naive bound. The required integral image introduces a $\Theta(N)$ memory footprint.

\begin{table}[h]
\centering
\caption{Asymptotic complexity of Classical Box-Counting.}
\label{tab:bc_complexity}
\begin{tabular}{lcc}
\hline
\textbf{Formulation} & \textbf{Time} & \textbf{Aux.\ Space} \\
\hline
Naive     & $\mathcal{O}(mN)$ & $\Theta(1)$ \\
Optimized & $\Theta(N)$       & $\Theta(N)$ \\
\hline
\end{tabular}
\end{table}

% ------------------------------------------------------------
\subsection{Gliding Box (GB)}
\label{subsec:gb}
% ------------------------------------------------------------

\subsubsection{Naive Formulation}

Algorithm~\ref{alg:gb_naive} implements the dense sliding window mass calculation defined in Eq.~\eqref{eq:gb_mass}. For a fixed scale $L$, scanning $L^2$ pixels across $\Theta(N)$ valid window positions yields a per-scale cost of:

\begin{equation}
  T_{\text{GB,naive}}(L)
    = \Theta(N) \cdot \Theta(L^{2})
    = \Theta(NL^{2}).
  \label{eq:gb_naive_scale}
\end{equation}

Summing over $m$ scales, the total execution time is heavily penalized by larger box sizes:

\begin{equation}
  T_{\text{GB,naive}}(N, m)
    = \Theta\!\left(N \sum_{L \in \mathcal{S}} L^{2}\right).
  \label{eq:gb_naive_total}
\end{equation}

The per-scale mass histogram occupies $\mathcal{O}(L_{\max}^{2})$ auxiliary space, which is generally negligible.

\begin{algorithm}[t]
\caption{Naive Gliding Box}
\label{alg:gb_naive}
\begin{algorithmic}[1]
\Require Binary image $I[0\ldots W{-}1][0\ldots H{-}1]$, scale set $\mathcal{S}$
\Ensure Pairs $\{(L,\, \mu(L),\, \Lambda(L)) : L \in \mathcal{S}\}$
\State Binarize $I$ via Otsu's threshold $t^{*}$ \cite{Otsu1979}
\For{each $L \in \mathcal{S}$}
  \State $H[0 \ldots L^{2}] \leftarrow 0$ \hfill\Comment{mass-frequency histogram}
  \For{$x \leftarrow 0$ \textbf{to} $W - L$}
    \For{$y \leftarrow 0$ \textbf{to} $H - L$}
      \State $\mathit{mass} \leftarrow 0$
      \For{$dx \leftarrow 0$ \textbf{to} $L - 1$}
        \For{$dy \leftarrow 0$ \textbf{to} $L - 1$}
          \State $\mathit{mass} \leftarrow \mathit{mass} + I[x{+}dx][y{+}dy]$
        \EndFor
      \EndFor
      \State $H[\mathit{mass}] \leftarrow H[\mathit{mass}] + 1$
    \EndFor
  \EndFor
  \State $T \leftarrow (W{-}L{+}1)(H{-}L{+}1)$
  \State Compute $\mu(L)$ and $\mu^{(2)}(L)$ from $H$ and $T$
  \State Compute $\Lambda(L)$ via Eq.~\eqref{eq:lacunarity}
  \State \textbf{yield} $(L,\, \mu(L),\, \Lambda(L))$
\EndFor
\State Estimate $D_B$ by least-squares regression on $\{(\log L_i,\, \log \mu(L_i))\}$
\end{algorithmic}
\end{algorithm}

% ............................................................
\subsubsection{Optimized Formulation}

Algorithm~\ref{alg:gb_opt} applies the integral-image optimization to the GB method. Pre-computing $P$ (Eq.~\eqref{eq:prefix}) reduces each dense window summation to $\mathcal{O}(1)$. If BC and GB operate in the same pipeline, this $\Theta(N)$ data structure is entirely shared.

\begin{algorithm}[t]
\caption{Optimized Gliding Box (Integral Image)}
\label{alg:gb_opt}
\begin{algorithmic}[1]
\Require Binary image $I[0\ldots W{-}1][0\ldots H{-}1]$, scale set $\mathcal{S}$, integral image $P$
\Ensure Pairs $\{(L,\, \mu(L),\, \Lambda(L)) : L \in \mathcal{S}\}$
\For{each $L \in \mathcal{S}$}
  \State $H[0 \ldots L^{2}] \leftarrow 0$
  \For{$x \leftarrow 0$ \textbf{to} $W - L$}
    \For{$y \leftarrow 0$ \textbf{to} $H - L$}
      \State $\mathit{mass} \leftarrow \sigma(x,\, y,\, x{+}L{-}1,\, y{+}L{-}1)$ \hfill\Comment{Eq.~\eqref{eq:rect_sum}, $\mathcal{O}(1)$}
      \State $H[\mathit{mass}] \leftarrow H[\mathit{mass}] + 1$
    \EndFor
  \EndFor
  \State $T \leftarrow (W{-}L{+}1)(H{-}L{+}1)$
  \State Compute $\mu(L)$ and $\mu^{(2)}(L)$ from $H$ and $T$
  \State Compute $\Lambda(L)$ via Eq.~\eqref{eq:lacunarity}
  \State \textbf{yield} $(L,\, \mu(L),\, \Lambda(L))$
\EndFor
\State Estimate $D_B$ by least-squares regression on $\{(\log L_i,\, \log \mu(L_i))\}$
\end{algorithmic}
\end{algorithm}

With $\mathcal{O}(1)$ mass evaluations, the per-scale complexity drops to the $\Theta(N)$ valid window positions:

\begin{equation}
  T_{\text{GB,opt}}(N, m)
    = \Theta(N) + \sum_{L \in \mathcal{S}} \Theta(N)
    = \Theta(mN).
  \label{eq:gb_opt_total}
\end{equation}

This represents a $\Theta(L^{2})$ acceleration per scale over the naive implementation. While BC visits $\Theta(N/\varepsilon^{2})$ disjoint cells, GB must visit $\Theta(N)$ overlapping positions regardless of scale. This algorithmic overhead enables GB to extract the richer mass-frequency distributions necessary for lacunarity estimation ($\Lambda$).

\begin{table}[h]
\centering
\caption{Asymptotic complexity of the Gliding Box method.}
\label{tab:gb_complexity}
\begin{tabular}{lcc}
\hline
\textbf{Formulation} & \textbf{Time} & \textbf{Aux.\ Space} \\
\hline
Naive     & $\Theta\!\left(N\!\sum_{L}L^{2}\right)$ & $\mathcal{O}(L_{\max}^{2})$ \\
Optimized & $\Theta(mN)$ & $\Theta(N)$ \\
\hline
\end{tabular}
\end{table}

\begin{table}[h]
\centering
\caption{BC vs.\ GB optimized formulations (per scale).}
\label{tab:bc_gb_comparison}
\resizebox{\columnwidth}{!}{%
\begin{tabular}{lccc}
\hline
\textbf{Method} & \textbf{Positions/scale} & \textbf{Cost/position} & \textbf{Per-scale cost} \\
\hline
BC (opt.) & $\Theta(N/\varepsilon^{2})$ & $\mathcal{O}(1)$ & $\Theta(N/\varepsilon^{2})$ \\
GB (opt.) & $\Theta(N)$ & $\mathcal{O}(1)$ & $\Theta(N)$ \\
\hline
\end{tabular}%
}
\end{table}

% ------------------------------------------------------------
\subsection{Differential Box-Counting (DBC)}
\label{subsec:dbc}
% ------------------------------------------------------------

\subsubsection{Naive Formulation}

Algorithm~\ref{alg:dbc_naive} directly implements the 3D grid partitioning detailed in Eq.~\eqref{eq:dbc_count}. Because it operates on grayscale data, no binarization is performed. Traversing the $\Theta(N/\varepsilon^2)$ cells requires an exhaustive $\Theta(\varepsilon^2)$ pixel scan per cell to determine the local intensity extrema $(I^{ij}_{\max}, I^{ij}_{\min})$.

\begin{algorithm}[t]
\caption{Naive Differential Box-Counting}
\label{alg:dbc_naive}
\begin{algorithmic}[1]
\Require Grayscale image $I[0\ldots W{-}1][0\ldots H{-}1]$, intensity range $G$, scale set $\mathcal{S}$
\Ensure Pairs $\{(\varepsilon,\, N(\varepsilon)) : \varepsilon \in \mathcal{S}\}$
\For{each $\varepsilon \in \mathcal{S}$}
  \State $\delta \leftarrow \lceil G \;/\; \lceil W/\varepsilon \rceil \rceil$
  \State $\mathit{total} \leftarrow 0$
  \For{$b_x \leftarrow 0$ \textbf{to} $\lceil W/\varepsilon \rceil - 1$}
    \For{$b_y \leftarrow 0$ \textbf{to} $\lceil H/\varepsilon \rceil - 1$}
      \State $I_{\max} \leftarrow -\infty$;\; $I_{\min} \leftarrow +\infty$
      \For{$dx \leftarrow 0$ \textbf{to} $\varepsilon - 1$}
        \For{$dy \leftarrow 0$ \textbf{to} $\varepsilon - 1$}
          \State $px \leftarrow b_x{\cdot}\varepsilon{+}dx$;\; $py \leftarrow b_y{\cdot}\varepsilon{+}dy$
          \If{$px < W \;\wedge\; py < H$}
            \State $I_{\max} \leftarrow \max(I_{\max},\, I[px][py])$
            \State $I_{\min} \leftarrow \min(I_{\min},\, I[px][py])$
          \EndIf
        \EndFor
      \EndFor
      \State $\mathit{total} \leftarrow \mathit{total} + \lceil (I_{\max}{-}I_{\min})/\delta \rceil + 1$
    \EndFor
  \EndFor
  \State \textbf{yield} $(\varepsilon,\, \mathit{total})$
\EndFor
\State Estimate $D_B$ by least-squares regression on $\{(\log \varepsilon_i,\, \log N(\varepsilon_i))\}$
\end{algorithmic}
\end{algorithm}

The resulting per-scale cost is $\Theta(N/\varepsilon^2) \cdot \Theta(\varepsilon^2) = \Theta(N)$, yielding a total time complexity of:

\begin{equation}
  T_{\text{DBC,naive}}(N, m) = \Theta(mN).
  \label{eq:dbc_naive_total}
\end{equation}

Although the asymptotic class mirrors naive BC, the hidden constant factor is strictly larger. Two intensity comparisons are required per pixel, and early-exit mechanisms are mathematically impossible since true extrema require scanning the entire cell. Auxiliary space remains $\Theta(1)$.

% ............................................................
\subsubsection{Optimized Formulation}

Algorithm~\ref{alg:dbc_opt} replaces the exhaustive $\Theta(\varepsilon^2)$ scan with offline 2D sparse tables~\cite{Liu2014}. The tables map regional minima and maxima queries to $\mathcal{O}(1)$ execution time.

\begin{algorithm}[t]
\caption{Optimized Differential Box-Counting (Sparse Tables)}
\label{alg:dbc_opt}
\begin{algorithmic}[1]
\Require Grayscale image $I[0\ldots W{-}1][0\ldots H{-}1]$, intensity range $G$, scale set $\mathcal{S}$
\Ensure Pairs $\{(\varepsilon,\, N(\varepsilon)) : \varepsilon \in \mathcal{S}\}$
\State Build 2D sparse table $\mathcal{T}_{\max}$ for range maximum \hfill \Comment{$\Theta(N \log N)$ time and space \cite{Liu2014}}
\State Build 2D sparse table $\mathcal{T}_{\min}$ for range minimum \hfill \Comment{$\Theta(N \log N)$ time and space}
\For{each $\varepsilon \in \mathcal{S}$}
  \State $\delta \leftarrow \lceil G \;/\; \lceil W/\varepsilon \rceil \rceil$
  \State $\mathit{total} \leftarrow 0$
  \For{$b_x \leftarrow 0$ \textbf{to} $\lceil W/\varepsilon \rceil - 1$}
    \For{$b_y \leftarrow 0$ \textbf{to} $\lceil H/\varepsilon \rceil - 1$}
      \State $x_1 \leftarrow b_x{\cdot}\varepsilon$;\; $y_1 \leftarrow b_y{\cdot}\varepsilon$
      \State $x_2 \leftarrow \min(x_1{+}\varepsilon{-}1,\,W{-}1)$;\; $y_2 \leftarrow \min(y_1{+}\varepsilon{-}1,\,H{-}1)$
      \State $I_{\max} \leftarrow \mathcal{T}_{\max}.\textsc{Query}(x_1, y_1, x_2, y_2)$ \hfill\Comment{$\mathcal{O}(1)$}
      \State $I_{\min} \leftarrow \mathcal{T}_{\min}.\textsc{Query}(x_1, y_1, x_2, y_2)$ \hfill\Comment{$\mathcal{O}(1)$}
      \State $\mathit{total} \leftarrow \mathit{total} + \lceil (I_{\max}{-}I_{\min})/\delta \rceil + 1$
    \EndFor
  \EndFor
  \State \textbf{yield} $(\varepsilon,\, \mathit{total})$
\EndFor
\State Estimate $D_B$ by least-squares regression on $\{(\log \varepsilon_i,\, \log N(\varepsilon_i))\}$
\end{algorithmic}
\end{algorithm}

Sparse table construction demands $\Theta(N \log N)$ time and space. Post-construction, traversing the $\Theta(N/\varepsilon^2)$ cells executes in unit time per cell. The total complexity is therefore bound by the preprocessing phase:

\begin{equation}
  T_{\text{DBC,opt}}(N, m)
    = \Theta(N \log N) + \sum_{\varepsilon \in \mathcal{S}} \Theta\!\left(\frac{N}{\varepsilon^{2}}\right)
    = \Theta(N \log N).
  \label{eq:dbc_opt_total}
\end{equation}

While computationally superior at large scales, DBC optimization introduces a severe $\Theta(N \log N)$ memory overhead, establishing a prominent time-memory tradeoff for grayscale analyses.

\begin{table}[h]
\centering
\caption{Asymptotic complexity of Differential Box-Counting.}
\label{tab:dbc_complexity}
\begin{tabular}{lcc}
\hline
\textbf{Formulation} & \textbf{Time} & \textbf{Aux.\ Space} \\
\hline
Naive     & $\Theta(mN)$        & $\Theta(1)$        \\
Optimized & $\Theta(N \log N)$  & $\Theta(N \log N)$ \\
\hline
\end{tabular}
\end{table}

% ------------------------------------------------------------
\subsection{Blanket Method (BM)}
\label{subsec:bm}
% ------------------------------------------------------------

\subsubsection{Naive Formulation}

Algorithm~\ref{alg:bm_naive} implements the sequential morphological erosion and dilation recurrences defined in Eqs.~\eqref{eq:blanket_u}--\eqref{eq:blanket_vol}. Each of the $\varepsilon_{\max}$ dilation levels requires examining the 4-connected neighborhood ($\mathcal{O}(1)$ work) for all $N$ pixels. 

\begin{algorithm}[t]
\caption{Naive Blanket Method}
\label{alg:bm_naive}
\begin{algorithmic}[1]
\Require Grayscale image $I[0\ldots W{-}1][0\ldots H{-}1]$, maximum dilation level $\varepsilon_{\max}$
\Ensure Pairs $\{(\varepsilon,\, A(\varepsilon)) : \varepsilon = 1,\ldots,\varepsilon_{\max}\}$
\State $u \leftarrow I$;\; $b \leftarrow I$ \hfill\Comment{initialize both blankets}
\State $V_{\text{prev}} \leftarrow \sum_{x,y} [u[x][y] - b[x][y]]$
\For{$\varepsilon \leftarrow 1$ \textbf{to} $\varepsilon_{\max}$}
  \State $u' \leftarrow$ new array of size $W \times H$
  \State $b' \leftarrow$ new array of size $W \times H$
  \For{$x \leftarrow 0$ \textbf{to} $W-1$}
    \For{$y \leftarrow 0$ \textbf{to} $H-1$}
      \State $u'[x][y] \leftarrow u[x][y] + 1$
      \State $b'[x][y] \leftarrow b[x][y] - 1$
      \For{each $(s,t) \in \mathcal{N}(x,y)$}
        \State $u'[x][y] \leftarrow \max(u'[x][y],\; u[s][t])$
        \State $b'[x][y] \leftarrow \min(b'[x][y],\; b[s][t])$
      \EndFor
    \EndFor
  \EndFor
  \State $u \leftarrow u'$;\; $b \leftarrow b'$
  \State $V \leftarrow \sum_{x,y}[u[x][y]-b[x][y]]$
  \State \textbf{yield} $\bigl(\varepsilon,\;(V - V_{\text{prev}})/2\bigr)$
  \State $V_{\text{prev}} \leftarrow V$
\EndFor
\State Estimate $D_B$ by least-squares regression on $\{(\log \varepsilon_i,\, \log A(\varepsilon_i))\}$
\end{algorithmic}
\end{algorithm}

The total execution time evaluates to:

\begin{equation}
  T_{\text{BM,naive}}(N, \varepsilon_{\max}) = \Theta(N \cdot \varepsilon_{\max}).
  \label{eq:bm_naive_total}
\end{equation}

Because blanket surfaces strictly depend on their prior states, parallelizing across scales is impossible. Auxiliary space requires $\Theta(N)$ to maintain the active and previous state buffers.

% ............................................................
\subsubsection{Optimized Formulation}

Evaluating morphological extrema across an expanded $(2\varepsilon+1) \times (2\varepsilon+1)$ window typically incurs an $\Theta(\varepsilon^2)$ penalty per pixel. Algorithm~\ref{alg:bm_opt} sidesteps this by exploiting the \textit{separability} of maximum and minimum filters. A 2D extrema window decomposes into successive 1D row and column passes. Using Lemire's monotonic deques \cite{Lemire2006}, each 1D pass resolves in $\Theta(N)$ amortized time.

\begin{algorithm}[t]
\caption{Optimized Blanket Method (Separable Monotonic Deque)}
\label{alg:bm_opt}
\begin{algorithmic}[1]
\Require Grayscale image $I[0\ldots W{-}1][0\ldots H{-}1]$, maximum dilation level $\varepsilon_{\max}$
\Ensure Pairs $\{(\varepsilon,\, A(\varepsilon)) : \varepsilon = 1,\ldots,\varepsilon_{\max}\}$
\State $u \leftarrow I$;\; $b \leftarrow I$
\State $V_{\text{prev}} \leftarrow \sum_{x,y} [u[x][y] - b[x][y]]$
\For{$\varepsilon \leftarrow 1$ \textbf{to} $\varepsilon_{\max}$}
  \State $u_{\text{inc}} \leftarrow u + 1$ \hfill\Comment{element-wise increment, $\Theta(N)$}
  \State $b_{\text{dec}} \leftarrow b - 1$
  \State $u' \leftarrow \textsc{SlidingMax2D}(u,\, K{=}3)$ \hfill\Comment{separable deque, $\Theta(N)$ \cite{Lemire2006}}
  \State $b' \leftarrow \textsc{SlidingMin2D}(b,\, K{=}3)$ \hfill\Comment{separable deque, $\Theta(N)$}
  \State $u \leftarrow \max(u_{\text{inc}},\, u')$ \hfill\Comment{element-wise, $\Theta(N)$}
  \State $b \leftarrow \min(b_{\text{dec}},\, b')$
  \State $V \leftarrow \sum_{x,y} [u[x][y] - b[x][y]]$
  \State \textbf{yield} $\bigl(\varepsilon,\; (V - V_{\text{prev}})/2\bigr)$
  \State $V_{\text{prev}} \leftarrow V$
\EndFor
\State Estimate $D_B$ by least-squares regression on $\{(\log \varepsilon_i,\, \log A(\varepsilon_i))\}$
\end{algorithmic}
\end{algorithm}

The total asymptotic time remains equivalent to the naive class:

\begin{equation}
  T_{\text{BM,opt}}(N, \varepsilon_{\max}) = \Theta(N \cdot \varepsilon_{\max}).
  \label{eq:bm_opt_total}
\end{equation}

However, the separable formulation guarantees that the computational constant remains flat regardless of the structuring element size. Table~\ref{tab:all_complexity} consolidates the bounds across all evaluated methods.

\begin{table}[h]
\centering
\caption{Asymptotic complexity of the Blanket Method.}
\label{tab:bm_complexity}
\begin{tabular}{lcc}
\hline
\textbf{Formulation} & \textbf{Time} & \textbf{Aux.\ Space} \\
\hline
Naive     & $\Theta(N\varepsilon_{\max})$ & $\Theta(N)$ \\
Optimized & $\Theta(N\varepsilon_{\max})$ & $\Theta(N)$ \\
\hline
\end{tabular}
\end{table}

\begin{table}[h]
\centering
\caption{Consolidated asymptotic complexity of all four methods (optimized formulations).}
\label{tab:all_complexity}
\begin{tabular}{lccc}
\hline
\textbf{Method} & \textbf{Input} & \textbf{Time} & \textbf{Aux.\ Space} \\
\hline
BC  (opt.) & Binary    & $\Theta(N)$           & $\Theta(N)$        \\
GB  (opt.) & Binary    & $\Theta(mN)$          & $\Theta(N)$        \\
DBC (opt.) & Grayscale & $\Theta(N \log N)$    & $\Theta(N \log N)$ \\
BM  (opt.) & Grayscale & $\Theta(N\varepsilon_{\max})$ & $\Theta(N)$ \\
\hline
\end{tabular}
\end{table}

% ============================================================
%  SECTION 4 — METHODOLOGY
% ============================================================
\section{Methodology}
\label{sec:methodology}

% ------------------------------------------------------------
\subsection{Materials}
\label{subsec:materials}
% ------------------------------------------------------------

The experimental evaluation is conducted on the UCSB Breast
Cancer Histology dataset, provided by the Center for
Bio-Image Informatics of the University of California, Santa
Barbara~\cite{Gelasca2008}. The dataset comprises 58
hematoxylin and eosin~(H\&E) stained histological images of
breast tissue at a fixed resolution of $896 \times 768$
pixels, partitioned into two classes: 32 benign and 26
malignant samples. The class imbalance ratio of
approximately 1.23:1 is mild enough to not require
resampling for the fractal dimension estimation task, since
$D$ is computed independently per image and no
classifier is trained on the extracted features.

This dataset was selected for three reasons directly relevant
to the algorithmic focus of this work. First, its fixed
resolution ensures that the pixel count $N = 896 \times 768
= 688{,}128$ is identical across all images, making
wall-clock time comparisons across methods and
formulations free of resolution confounds. Second, the two
well-defined tissue classes — benign and malignant — provide
a concrete discriminatory benchmark against which the
effectiveness of each algorithm's $D$ estimate can be
assessed, following the methodology established in prior
fractal-based histological classification
work~\cite{Roberto2021,Borgue2025}. Third, the dataset is
publicly available as a Kaggle dataset
(\texttt{ucsb-breast-cancer-histology}), enabling direct
reproducibility of all experiments within the execution
environment described in Section~\ref{subsec:environment}.

All images are loaded as 8-bit grayscale arrays ($G = 256$)
prior to processing. For the binary-domain methods~(BC and
GB), a binarisation step is applied using Otsu's global
threshold~\cite{Otsu1979} before the fractal dimension
estimation loop; this step is timed and reported separately
from the algorithmic core, as described in
Section~\ref{subsec:environment}.

% ------------------------------------------------------------
\subsection{Experimental Environment and Protocol}
\label{subsec:environment}
% ------------------------------------------------------------

All experiments are executed on Kaggle Notebooks
(\url{https://www.kaggle.com}) using a CPU-only session,
to ensure that the measured execution times reflect the
sequential computational complexity of each algorithm
without GPU acceleration. The hardware configuration
provided by the Kaggle CPU session consists of an Intel
Xeon processor with two virtual cores and approximately
13~GB of available RAM. The software environment is fixed
at Python~3.10 with NumPy~1.26, SciPy~1.11, and
scikit-image~0.22; the complete \texttt{requirements.txt} is included in the
\textit{github} repository to ensure reproducibility.
\footnote{Source code available at:
\href{https://https://github.com/gabrielmbrum} {\texttt{https://github.com/gabrielmbrum}}.}

The scale set used across all four methods is fixed as
$\mathcal{S} = \{2, 4, 8, 16, 32, 64\}$, comprising $m = 6$
scales spanning approximately 1.5 decades of spatial
frequency. This range respects the lower bound imposed by
the pixel size and the upper bound of $L_{\max} = 64 \ll
\sqrt{N} \approx 830$, ensuring that the covering assumption
underlying each algorithm remains valid throughout the scale
range. The same $\mathcal{S}$ is used for all methods to
guarantee that differences in execution time and $R^{2}$
are attributable solely to algorithmic structure rather
than to differences in scale coverage.

Each algorithm — in both its naive and optimised
formulations — is executed 10 times per image. The first
run is discarded as a warm-up to eliminate cold-cache and
interpreter overhead effects; the remaining 9 runs are used
to compute the mean wall-clock time $\bar{t}$ and standard
deviation $\sigma_t$. Wall-clock time is measured using
Python's \texttt{time.perf\_counter()}, which provides
sub-microsecond resolution on the target platform. Runs
exhibiting a coefficient of variation

\begin{equation}
  CV = \frac{\sigma_t}{\bar{t}} \times 100\%
  \label{eq:cv}
\end{equation}

\noindent exceeding 5\% are flagged and re-evaluated with
additional runs until $CV < 5\%$, following the adaptive
stability criterion recommended for wall-clock
benchmarking~\cite{Flemming1986}. This threshold ensures
that the reported means are statistically stable without
imposing an unnecessarily large number of repetitions.

Pre-processing time is measured and reported separately
from the algorithmic core. Concretely, for BC and GB, the
wall-clock time of the Otsu binarisation step is recorded
independently of the subsequent coverage computation,
allowing the contribution of the pre-processing stage to
the total execution budget to be quantified and discussed
in Section~\ref{sec:results}. For DBC and BM, which
operate directly on the grayscale input, no binarisation
is performed and pre-processing time is zero by definition.

% ------------------------------------------------------------
\subsection{Evaluation Metrics}
\label{subsec:metrics}
% ------------------------------------------------------------

Three orthogonal evaluation axes are considered, each
capturing a distinct dimension of the algorithmic
trade-off under study.

\subsubsection{Execution Time}

The primary time metric is the mean wall-clock time
$\bar{t}$ (in seconds) averaged over all 58 images and
9 valid runs per image. To assess the agreement between
theoretical predictions and empirical scaling behaviour,
$\bar{t}$ is plotted as a function of image size $N$ across
sub-images of increasing resolution, and the empirical
slope of the $\log \bar{t}$ versus $\log N$ curve is
compared against the theoretical exponent derived in
Section~\ref{sec:algorithms}. Deviations between empirical
and theoretical slopes are attributed to constant factors,
memory hierarchy effects, and interpreter overhead, and are
discussed qualitatively in Section~\ref{sec:results}.

\subsubsection{Memory Consumption}

Auxiliary memory consumption is measured as the peak
heap allocation incurred by each algorithm beyond the
read-only input buffer, using Python's
\texttt{tracemalloc} module. The input buffer of size
$\Theta(N)$ bytes is excluded from the measurement by
recording the allocation differential between the start
and end of the algorithmic core. This isolates the
auxiliary space — the integral image for BC and GB, the
sparse tables for DBC, and the blanket arrays for BM —
and allows direct comparison against the $\Theta(N)$ and
$\Theta(N \log N)$ theoretical bounds derived in
Section~\ref{sec:algorithms}.

\subsubsection{Effectiveness}

Effectiveness is defined as the capacity of each
algorithm's $D$ estimate to statistically discriminate
between the benign and malignant tissue classes of the
UCSB dataset. Formally, for each method, $D$ is computed
for every image and the class-conditional distributions
$\{D_i : i \in \text{benign}\}$ and $\{D_j : j \in
\text{malignant}\}$ are obtained. The primary effectiveness
metric is the inter-class separation

\begin{equation}
  \Delta D
    = \bar{D}_{\text{malignant}} - \bar{D}_{\text{benign}},
  \label{eq:delta_d}
\end{equation}

\noindent where $\bar{D}_{\text{malignant}}$ and
$\bar{D}_{\text{benign}}$ denote the respective class
means. A larger $|\Delta D|$ indicates that the method
assigns meaningfully different fractal dimensions to the
two tissue types, reflecting the clinically established
observation that malignant tissues exhibit higher surface
roughness and structural complexity than benign
ones~\cite{Lopes2011,Rangayyan2007}. Statistical
significance of $\Delta D$ is assessed via Welch's
two-sample $t$-test at significance level $\alpha = 0.05$,
which does not assume equal variances across classes and
is appropriate given the mild class imbalance of the
dataset~\cite{Welch1947}.

Two complementary effectiveness indicators are also
reported. The coefficient of determination $R^{2}$ of the
log-log regression used to estimate $D$ from each image
quantifies the degree to which the object exhibits
consistent power-law scaling across the chosen scale set;
values close to unity are a necessary condition for
interpreting $D$ as a valid fractal descriptor~\cite{
Lopes2011}. Additionally, for the binary methods~(BC and
GB), the sensitivity of $\Delta D$ to the Otsu threshold
is evaluated by perturbing $t^{*}$ by $\pm 10\%$ and
measuring the resulting variation in $D$, providing a
quantitative characterisation of the threshold sensitivity
trade-off that distinguishes binary from grayscale methods.

% ============================================================
%  SECTION 5 — RESULTS AND DISCUSSION
% ============================================================
\section{Results and Discussion}
\label{sec:results}

% ------------------------------------------------------------
\subsection{Pre-processing Cost}
\label{subsec:preproc}
% ------------------------------------------------------------

The binarization step required by the BC and GB methods, performed via Otsu's global threshold~\cite{Otsu1979}, was timed independently on all 58 images. The mean binarization time was $1.78 \pm 0.60$~ms per image. This represents less than $1\%$ of the total execution budget for either method in their optimized formulations. Consequently, this pre-processing stage is practically negligible and does not constitute a computational bottleneck for clinical deployment. This isolation confirms that the execution times reported subsequently for BC and GB accurately reflect the fundamental cost of their algorithmic cores.

% ------------------------------------------------------------
\subsection{Execution Time: Empirical vs. Theoretical}
\label{subsec:res_time}
% ------------------------------------------------------------

Table~\ref{tab:res_time} summarizes the mean wall-clock execution times and empirical speedups.

\begin{table}[h]
\centering
\caption{Mean wall-clock execution time per image. Speedup $= \bar{t}_{\text{naive}} / \bar{t}_{\text{opt}}$.}
\label{tab:res_time}
\resizebox{\columnwidth}{!}{%
\begin{tabular}{lrrr}
\hline
\textbf{Method} & \textbf{Naive (ms)} & \textbf{Optimised (ms)} & \textbf{Speedup} \\
\hline
BC  & $186.0 \pm 51.0$  & $248.4 \pm 3.4$   & $0.75\times$ \\
GB  & $2718.5 \pm 33.0$ & $45.8 \pm 2.2$    & $59.3\times$ \\
DBC & $1027.1 \pm 19.2$ & $119.5 \pm 43.0$  & $8.6\times$  \\
BM  & $409.7 \pm 23.6$  & $306.8 \pm 20.1$  & $1.3\times$  \\
\hline
\end{tabular}
}
\end{table}

The empirical execution times yield a central scientific finding of this study: \textit{asymptotic superiority alone is insufficient to predict practical performance.} This is starkly evidenced by the BC method, where the structurally optimized formulation ($248.4$~ms) executed slower than the naive approach ($186.0$~ms), yielding a paradoxical speedup of $0.75\times$. Although the integral image mathematically reduces the $\mathcal{O}(N)$ worst-case per-scale traversal to a strict $\Theta(N/\varepsilon^2)$, practical performance is dictated by memory hierarchy effects and implementation-level characteristics. The integral image introduces a $\approx 2.6$~MB \texttt{int32} array that exceeds the L2 cache of the evaluation environment, inducing substantial memory-access latency. Conversely, the naive formulation accesses the original binary array in a strictly sequential pattern, maximizing hardware prefetcher efficiency. Furthermore, histological images possess dense foreground regions, allowing the naive inner loop to trigger early-exit conditions almost immediately, reducing the effective constant factor well below theoretical worst-case predictions. 

In contrast, theoretical improvements aligned with massive practical gains for the GB method, achieving a $59.3\times$ speedup ($45.8$~ms vs $2718.5$~ms). Because the naive GB evaluates $\Theta(N)$ overlapping windows per scale without early-exits, the $\Theta(L^2)$ pixel additions strictly manifest in practice (costing $\approx 2.8 \times 10^9$ operations at $L=64$). The integral image optimization completely eliminates this $\Theta(L^2)$ bottleneck, proving indispensable for the clinical deployment of GB at high resolutions.

The optimized DBC achieved an $8.6\times$ speedup ($119.5$~ms vs $1027.1$~ms) by utilizing sparse tables to eliminate the exhaustive $\Theta(\varepsilon^2)$ extrema search. Since true extrema computation permits no early-exit, removing the scan yields immediate wall-clock benefits. However, the high standard deviation ($\pm 43.0$~ms) stems directly from operating-system memory jitter during the massive $\Theta(N \log N)$ sparse table allocation.

Finally, the BM formulation yielded a modest $1.3\times$ speedup ($306.8$~ms vs $409.7$~ms). As derived theoretically, both formulations share a $\Theta(N \varepsilon_{\max})$ complexity. The minor gain arises from vectorized NumPy filtering, which minimizes Python interpreter overhead compared to explicit neighborhood loops. Substantial speedups for BM are only anticipated when expanding beyond the $3 \times 3$ morphological structuring element.

% ------------------------------------------------------------
\subsection{Memory Consumption}
\label{subsec:res_mem}
% ------------------------------------------------------------

\begin{table}[h]
\centering
\caption{Mean peak auxiliary memory per image.}
\label{tab:res_mem}
\begin{tabular}{lrr}
\hline
\textbf{Method} & \textbf{Naive (KB)} & \textbf{Optimised (KB)} \\
\hline
BC  & $0.4$       & $5{,}376.6$ \\
GB  & $0.3$       & $5{,}376.6$ \\
DBC & $0.4$       & $413{,}954.6$ \\
BM  & $1.4$       & $1.3$ \\
\hline
\end{tabular}
\end{table}

Table~\ref{tab:res_mem} highlights memory as a first-class design consideration. While naive formulations maintain $\mathcal{O}(1)$ or minimal $\Theta(N)$ auxiliary space (operating effectively in-place), optimizations impose significant memory-performance trade-offs.

The BC and GB optimized variants require $\approx 5.25$~MB to store the shared integral image, a negligible overhead for modern desktop systems. However, the DBC optimized formulation introduces an exceptionally severe $\approx 404$~MB memory footprint to maintain the 2D sparse tables. While theoretically constrained to $\Theta(N \log N)$, the absolute magnitude of this requirement has profound practical implications. Allocating nearly half a gigabyte of auxiliary memory for a single $896 \times 768$ image poses strict deployment limitations for embedded diagnostic systems, parallelized high-throughput clinical cloud environments, or application to gigapixel Whole-Slide Images (WSIs). In such environments, practitioners may be forced to revert to the slower naive DBC implementation to satisfy memory constraints.

The BM approach proves highly advantageous in this regard, operating efficiently with minimal $\Theta(N)$ auxiliary arrays ($\approx 1.3$~KB peak overhead), making it the most memory-efficient grayscale algorithm evaluated.

% ------------------------------------------------------------
\subsection{Effectiveness: Fractal Dimension and Discriminatory Power}
\label{subsec:res_effectiveness}
% ------------------------------------------------------------

\begin{table}[h]
\centering
\caption{Mean fractal dimension $\bar{D}$ per class, inter-class separation $\Delta D$, mean $R^{2}$, and Welch $t$-test result ($\alpha = 0.05$). Naive and optimized formulations produced identical $D$ values.}
\label{tab:res_effectiveness}
\resizebox{\columnwidth}{!}{%
\begin{tabular}{lcccccc}
\hline
\textbf{Method} & $\bar{D}_{\text{ben}}$ & $\bar{D}_{\text{mal}}$ & $\Delta D$ & $\bar{R}^{2}$ & $p$-value & \textbf{Sig.} \\
\hline
BC  & $1.877 \pm 0.052$ & $1.790 \pm 0.086$ & $-0.087$ & $0.999$ & $0.0001$ & \checkmark \\
GB  & $-1.995 \pm 0.005$ & $-1.992 \pm 0.009$ & $+0.003$ & $1.000$ & $0.149$  & $\times$ \\
DBC & $2.112 \pm 0.038$ & $2.135 \pm 0.036$ & $+0.024$ & $0.992$ & $0.020$  & \checkmark \\
BM  & $0.535 \pm 0.087$ & $0.609 \pm 0.070$ & $+0.074$ & $0.987$ & $0.001$  & \checkmark \\
\hline
\end{tabular}
}
\end{table}

Table~\ref{tab:res_effectiveness} reports the statistical separation capabilities of the estimators. Under the evaluated dataset and experimental conditions, BC and BM exhibit the most robust discriminatory power. BC yields the largest absolute separation ($|\Delta D| = 0.087$, $p = 0.0001$). Malignant tissues manifest lower fractal dimensions in the binary domain, reflecting highly fragmented, less space-filling structural patterns compared to the well-organized glandular formations of benign tissues. Similarly, BM accurately captures the heightened structural irregularity of the malignant grayscale intensity surface, yielding strongly significant separation ($p = 0.001$). 

Conversely, the GB method fails to discriminate classes using the fractal dimension ($p = 0.149$). Many readers may interpret the negative $\bar{D}$ values ($\approx -1.99$) as computational errors; however, these strictly result from the adopted formulation and regression convention. Because GB fits the log-log regression to the \textit{mean box mass} $\mu(L)$—which scales proportionally to $L^2$ in uniform regions—rather than a traditional covering count, the slope converges to $+2$. Consequently, the fractal dimension defined as $-\text{slope}$ converges to $-2$. The failure of GB to discriminate via $D$ demonstrates that mean box mass carries limited structural signal. Instead, its discriminatory utility relies entirely on lacunarity, which exhibited strong inter-class separation ($\bar{\Lambda}_{\text{ben}} = 0.421$ vs $\bar{\Lambda}_{\text{mal}} = 0.832$). This confirms that the spatial heterogeneity (clustering) of the tumor cells, rather than their mean mass, drives histological classification.

Finally, while BC provides strong statistical separation, a robustness analysis revealed a mean $D$ shift of $0.085$ when perturbing the Otsu threshold by $\pm 10\%$. Because this variation is essentially equivalent to the inter-class separation ($|\Delta D| = 0.087$), improper pre-processing can entirely neutralize the discriminatory signal, highlighting a severe vulnerability of binary algorithms.

% ------------------------------------------------------------
\subsection{Trade-off Synthesis}
\label{subsec:tradeoffs}
% ------------------------------------------------------------

\begin{table}[h]
\centering
\caption{Consolidated trade-off summary across all optimized formulations. $\checkmark$ = best in class; $\sim$ = acceptable; $\times$ = weakest.}
\label{tab:tradeoffs}
\resizebox{\columnwidth}{!}{%
\begin{tabular}{lcccc}
\hline
\textbf{Criterion} & \textbf{BC} & \textbf{GB} & \textbf{DBC} & \textbf{BM} \\
\hline
Execution time      & $\checkmark$ & $\checkmark$ & $\sim$ & $\sim$ \\
Auxiliary memory    & $\sim$       & $\sim$       & $\times$ & $\checkmark$ \\
$R^{2}$ quality     & $\checkmark$ & $\checkmark$ & $\sim$ & $\sim$ \\
Discriminatory power ($|\Delta D|$) & $\checkmark$ & $\times$ & $\sim$ & $\checkmark$ \\
Threshold robustness & $\times$   & $\times$     & $\checkmark$ & $\checkmark$ \\
\hline
\end{tabular}
}
\end{table}

Table~\ref{tab:tradeoffs} synthesizes the fundamental trade-offs guiding algorithm selection. The binary-versus-grayscale dichotomy is paramount: binary methods (BC, GB) execute rapidly but suffer from extreme threshold sensitivity, potentially destroying their discriminatory validity if staining protocols vary. Grayscale methods (DBC, BM) bypass binarization, guaranteeing superior reproducibility across clinical cohorts at the expense of higher computational complexity and memory utilization. 

Practically, if strict image standardization is guaranteed and speed is critical, the naive BC is highly recommended. If threshold robustness is required and memory resources are constrained, BM provides strong, statistically significant discrimination without the crippling memory overhead of DBC. The optimized GB should be strictly reserved for pipelines requiring lacunarity evaluation.

% ============================================================
%  SECTION 7 — CONCLUSION
% ============================================================
\section{Conclusion}
\label{sec:conclusion}

This study demonstrates a fundamental reality of algorithm engineering in medical image analysis: asymptotic complexity alone is insufficient to predict practical performance. While structural optimizations—integral images, sparse tables, and separable monotonic deques—successfully reduced theoretical time bounds, their real-world deployment exposed critical bottlenecks dictated by memory hierarchy, cache utilization, and data-dependent early-exit conditions.

The empirical evidence underscores severe time-memory trade-offs across the evaluated estimators. Optimizing the Gliding Box (GB) method proved indispensable, yielding a 59.3$\times$ speedup that transformed a computationally intractable formulation into a clinically deployable solution. Similarly, the Differential Box-Counting (DBC) method achieved an 8.6$\times$ acceleration, but introduced a massive $\approx$404 MB memory penalty that limits its applicability in constrained environments. In stark contrast, the structurally optimized Classical Box-Counting (BC) method was paradoxically slower ($0.75\times$) than its naive counterpart, demonstrating that the overhead of large auxiliary data structures can easily eclipse theoretical gains when cache misses and the loss of early-exit conditions dominate execution. The Blanket Method (BM) maintained an exceptionally low memory footprint but yielded only a marginal $1.3\times$ speedup under standard morphological configurations.

Beyond computational efficiency, the findings highlight a critical dichotomy between binary and grayscale methods regarding discriminatory effectiveness. BC and BM provided the strongest statistical separation between benign and malignant histological tissues ($p < 0.01$). However, binary algorithms (BC, GB) exhibited severe vulnerability to pre-processing artifacts; a mere $\pm10\%$ perturbation in the binarization threshold was sufficient to completely neutralize the diagnostic signal of BC. Grayscale methods (DBC, BM) bypass this vulnerability entirely, trading higher computational demand for structural robustness. Furthermore, the analysis revealed that GB cannot reliably discriminate tissues using the fractal dimension; its diagnostic utility relies exclusively on lacunarity, capturing spatial heterogeneity rather than mean structural density.

Based on these consolidated findings, we provide the following practical recommendations for algorithm selection in medical image pipelines:
\begin{itemize}
    \item \textbf{Classical Box-Counting (BC):} Recommended strictly when rapid execution is critical and image binarization protocols are perfectly standardized to prevent threshold-induced signal degradation. The naive formulation is preferred for standard-resolution images to maximize cache efficiency.
    \item \textbf{Blanket Method (BM):} The optimal choice when threshold robustness and strong discriminatory power are required under strict memory constraints.
    \item \textbf{Differential Box-Counting (DBC):} A robust grayscale alternative that avoids binarization artifacts, provided the deployment environment can comfortably accommodate massive $\Theta(N \log N)$ memory allocations.
    \item \textbf{Gliding Box (GB):} Should be explicitly reserved for analytical pipelines that require the quantification of spatial clustering and heterogeneity via lacunarity, as its fractal dimension estimate lacks discriminatory power.
\end{itemize}

Ultimately, the design of future medical image analysis pipelines must treat hardware constraints—particularly peak memory allocation and cache behavior—as primary design variables alongside diagnostic power. Future investigations should evaluate these implementations across higher-resolution modalities (e.g., gigapixel Whole-Slide Images) and parallelized architectures (GPU/SIMD) to further stress-test the time-memory-accuracy boundaries established in this work.

% a bibliografia deve estar no arquivo references.bib em formato bibtex
\renewcommand{\refname}{References}
\bibliography{references}

\end{document}